\begin{afterword}
\summaryhead

The post argues that many causes gain when people become more rational, because their case takes a little more rationality than usual to see. Such causes could share standard material on rationality and accept that they will not capture all the value they create, since they will also capture some of what others create.

For this to work, each cause must be willing to stop talking about itself, and not treat supporters as a fixed supply. An organization's message ``This is the best thing you can do'' had fared worse than the X Prize Foundation's ``Here is a cool thing you can do.'' People are not expected-utility maximizers but ``altruistic akrasics,'' the author argues, so claiming to be best raises the burden of proof and invites comparisons that sap motivation; arguing over the best cause may reduce the total supply of altruism. Dollars are not fungible in practice: they come from different mental accounts and cost different amounts of willpower.

So causes should say ``Here is a cool rationalist project'' in general forums, and leave claims of top priority to specialized venues. Projects that are not real tasks for humanity are excluded. Rationality itself is not the most important thing, only ``a cool thing'' many people have in common.

\responsehead

The post is a proposal for cooperation among causes, and its psychological claims are hedged (``may,'' ``perhaps''). Its central empirical premise, that giving is not a fixed pie, has since been tested in the field, and the results mostly support it. After deadly tornadoes, households gave more without giving less to other local charities; matching grants on a charity website did not crowd out gifts to similar requests; one study of major appeals found that they raise total giving while shifting the timing of gifts to other charities. A direct-mail experiment found that adding scientific evidence of a charity's impact did not change giving overall, which fits, indirectly, the claim that ``best'' adds little motivation to ``cool.''

The evidence the post itself offers for its main rule is thinner. It is one comparison, reported secondhand, between an unnamed organization and the X Prize Foundation. The two differ in more than their slogans: the X Prize Foundation offers prizes for public, set goals, so one comparison cannot separate the effect of the message. The next paragraph supplies reasons, but no test.

In short: a hedged proposal for cooperation among causes, whose rule against claiming to be ``the best'' rests on one secondhand comparison and the reasons that follow it.
\end{afterword}
